%==============================================================================
% related_work.tex — Related Work for PP-HFM
% Self-contained \section{}; \input this into main.tex (never compile alone).
%
% Structured after the house related-work template: one dense paragraph per
% subsection, each closing on the limitation that motivates the next, and a
% final Summary and Research Gap that enumerates the observations before
% stating how the present work addresses them.
%
% POSITIONING RULE (non-negotiable, see docs/roadmap.md):
%   Claim: "we quantify what has not been quantified, and audit it properly."
%   NOT:   "we invented patient-level DP."  <- desk reject
% User-level DP (McMahan 2018), parameter-efficient DP fine-tuning (Li 2021,
% Yu 2021) and DP quantile release (Smith 2011) are PRIOR ART and are named as
% such before any gap is claimed.
%==============================================================================

\section{Related Work}
\label{sec:related}

\subsection{Federated Learning and Hierarchical Aggregation}

The formulation of distributed training as the communication of model updates
rather than of data \cite{ref1} established the basis upon which
subsequent federated systems were constructed. In its flat form, however, every
participating client transmits to a central server in every round, so that the
wide-area link rather than local computation determines how frequently the
global model may be revised. Client--edge--cloud designs \cite{ref2}
addressed this by interposing an intermediate tier at which several rounds of
local optimisation are absorbed before one round of central aggregation is
performed. HED-FL \cite{ref3} developed the principle for
constrained deployments, clustering devices dynamically at the edge and
demonstrating that the avoidance of long-range transmission substantially
extends the operational lifetime of an installation, while HCEF
\cite{ref4} introduced heterogeneity awareness, adjusting per-device
computation frequency and compression ratio online so that the slowest
participant does not dictate the duration of a round. Mhaisen \emph{et al.}
\cite{ref5} treated the assignment of clients to aggregators as
the object of optimisation, showing that an assignment informed by the
statistical properties of local data reduces model divergence under
non-identically distributed conditions. These works establish the
communication and energy case for hierarchy convincingly, and the present study
adopts rather than re-derives it. None of them, however, carries a formal
privacy guarantee: HED-FL provides none, and the assignment strategy of
\cite{ref5} requires clients to disclose local label
statistics, which is itself a release of information about the protected
population and is difficult to reconcile with a strict budget.

\subsection{Differential Privacy in Federated Learning}

Differentially private stochastic gradient descent \cite{ref6}
supplies the guarantee that hierarchical designs omit, bounding the influence of
a training record by clipping its gradient to a fixed norm and perturbing the
aggregate with calibrated noise, with the accumulated cost tracked by a moments
or R\'enyi accountant \cite{ref7} within the framework of
\cite{ref8}. Within federated learning the mechanism has been
developed in two principal directions. FL2DP \cite{ref9} combines an
exponential noise mechanism with a shuffling entity that severs the association
between a client and the update it contributes, reporting improved accuracy at
an equivalent formal guarantee, while RoPPFL \cite{ref10} pairs local
differential privacy with similarity-based robust aggregation performed at the
edge, so that the intermediate tier serves as a verifier against corrupted
updates as well as a relay. Both establish that a formal guarantee is
attainable in the federated setting at acceptable cost. Both, however, are
organised flatly with respect to the property considered here: FL2DP provides
differential privacy without a hierarchy in the sense that the edge tier plays
no part in the accounting, and neither examines how the interposition of such a
tier alters the frequency with which each client's budget is expended, which is
the quantity upon which the calibration of the noise ultimately depends.

\subsection{The Unit to Which a Guarantee Attaches}

The distinction between a guarantee stated over records and one stated over
individuals is long established and is not a contribution of the present work.
McMahan \emph{et al.} \cite{ref11} introduced user-level
differential privacy for federated language models, bounding the contribution of
each user rather than of each example, and the group-privacy relation by which
one may be converted into the other is standard \cite{ref8}.
The mechanism is therefore available; what is not available is a measurement of
what it costs in the setting that most requires it. In clinical monitoring the
grouping of records under an individual is not a property of a few atypical
participants but of every participant, since a single episode of care yields
many time-aligned records, and the conversion factor is consequently applied to
the whole population rather than to its tail. The literature supplies neither
the utility forgone when the unit is moved from the record to the individual on
genuinely paired multimodal clinical data, nor the manner in which that cost
varies with the number of records an individual contributes, and in the absence
of the latter a single measurement cannot be distinguished from a property of
the particular cohort on which it was obtained.

\subsection{Parameter-Efficient Private Adaptation}

The perturbation that a parameter vector must absorb under a fixed budget grows
with its dimensionality, so that private training of an entire network degrades
severely at the budgets a clinical claim requires. Restricting the private
optimisation to a compact trainable subset situated above frozen pretrained
representations is an established remedy, and Li \emph{et al.}
\cite{ref12} and Yu \emph{et al.} \cite{ref13} demonstrate
that parameter-efficient differentially private fine-tuning is markedly more
robust at small $\varepsilon$ than end-to-end private training. The present
work adopts this recipe and claims no novelty for it. The recipe nevertheless
carries a precondition that is frequently left implicit: the corpus upon which
the frozen representations are obtained must be disjoint from the protected
population, since otherwise the representations have already been derived from
that population without a guarantee and the budget subsequently reported for
the trainable component does not describe the system as a whole. Where the
precondition is honoured it is seldom priced, so the utility conceded in order
to obtain a guarantee that is valid end-to-end remains unquantified; and the
consequences of the recipe for what the resulting model retains of its training
population have received little attention.

\subsection{Membership Inference as Empirical Auditing}

A differential privacy guarantee is an upper bound upon the influence of a
participant and does not by itself establish what an adversary attains in
practice, for which membership inference supplies the complementary
measurement. Shokri \emph{et al.} \cite{ref14} introduced the
shadow-model formulation, and thresholding of the loss or confidence
\cite{ref15} remains the most commonly reported attack. Such
attacks are weak by construction: a single global threshold is confounded by
the varying difficulty of individual examples, so that an easily fitted
non-member may be scored above a poorly fitted member and the reported strength
constitutes a lower bound upon leakage rather than an estimate of it. Watson
\emph{et al.} \cite{ref16} demonstrated that calibrating against
per-example difficulty alters conclusions materially, and Carlini \emph{et al.}
\cite{ref17} formalised the calibration as a likelihood-ratio test
between per-example in- and out-distributions estimated from shadow models,
arguing further that attacks should be assessed by their true-positive rate at
low false-positive rate rather than by an average-case summary that conceals
confident identification of a minority. Two limitations persist. Membership is
almost invariably evaluated at the granularity of the training record, which is
precisely the quantity that a record-level guarantee already protects and not
the quantity a clinical guarantee concerns; and negative results are rarely
accompanied by evidence that the attack employed was capable of a positive one,
so that an absence of detected leakage cannot be distinguished from an
inadequate adversary.

\subsection{Multimodal and Clinical Federated Learning}

The motivation for federated learning is particularly clear in medicine, where
governance rather than volume prevents the assembly of data at a single site
\cite{ref18, ref19}. The majority of federated work,
including the systems discussed above \cite{ref9, ref3, ref4, ref10}, nevertheless treats the local objective as the
aggregation of generic parameter vectors and does not represent multiple
modalities per record, notwithstanding that clinical evidence is characteristically
multimodal and that the modalities available at one institution may differ from
those available at another. Where multimodal federated benchmarks have been
assembled, the modalities are frequently associated by a construction rule
rather than by common provenance, so that an apparent benefit of fusion may
reflect the regularity of the rule rather than genuinely complementary
information; a benchmark of this kind cannot support a claim about fusion at
all. Corpora in which both modalities originate from the same individual during
the same episode and are aligned at the point of capture
\cite{ref20}, or are synchronised in hardware
\cite{ref21}, avoid the difficulty entirely, since the association is a
property of the source and requires no defence.

\subsection{Uncertainty Quantification under a Privacy Budget}

Conformal prediction provides prediction sets carrying a distribution-free
coverage guarantee and requires neither retraining nor distributional
assumptions \cite{ref22, ref23}, which
recommends it for clinical deployment, where a point prediction unaccompanied
by a statement of confidence is of limited use. Its application within a
private system introduces a complication that is easily overlooked: the
procedure requires the release of a quantile of nonconformity scores computed
upon held-out data belonging to the protected population, and such a release is
a disclosure in exactly the sense the surrounding framework exists to bound.
Differentially private release of a quantile by the exponential mechanism is
itself standard \cite{ref24}, so the mechanism is not at issue.
What is at issue is the accounting: the cost of calibration composes with the
cost of training, and a system that reports only the latter understates the
budget it has expended, by a factor that grows with the number of occasions on
which calibration is performed.

\subsection{Real-Time Inference and Quality of Service at the Edge}

A federated system is trained once and queried continuously, and the constraint
that governs its clinical usefulness after deployment is the delay between a
measurement arriving and a decision being available, rather than the time taken
to reach convergence. The federated literature has addressed the latter
extensively and the former hardly at all: HCEF \cite{ref4} adapts computation
frequency and compression ratio online, but the quantity it optimises is the
duration of a training round; RoPPFL \cite{ref10} places verification at the
edge in order to reject corrupted updates, prioritising robustness over
response time; and HED-FL \cite{ref3} is concerned with the energy consumed in
reaching a model rather than the latency of using one. Where inference cost has
been considered it has generally been within the edge-computing literature
rather than the federated one \cite{ref25}, and the principal instrument is
reduced-precision arithmetic, integer quantisation in particular being well
established as a means of lowering both memory footprint and inference time on
commodity hardware \cite{ref26}. The consequence is a gap that matters
precisely for the systems this work targets: a federated design may satisfy a
strict privacy budget and still be unusable at the bedside if the delay between
observation and prediction exceeds the window in which a clinician can act,
and no existing federated system reports an end-to-end inference latency
measured under the privacy mechanism it deploys, alongside the communication
and energy costs incurred in producing the model that serves it.

\subsection{Summary and Research Gap}

These lines of work converge upon the setting from which the present study
proceeds, and five observations remain. First, hierarchy and formal privacy
have been developed largely in separation --- the hierarchical systems provide
no guarantee, and the private systems do not examine how an intermediate tier
alters the expenditure of the budget --- so that the interaction between the two
is unmeasured although each is individually well understood. Second, the
mechanism for stating a guarantee over individuals rather than records is
established, but its cost on genuinely paired clinical multimodal data is not,
and in the absence of any characterisation of how that cost varies with the
number of records an individual contributes, a single measurement cannot be
distinguished from an artefact of the cohort upon which it was obtained. Third,
parameter-efficient private adaptation is widely adopted, yet the disjointness
of the pretraining corpus upon which its validity depends is frequently left
implicit and, where honoured, is not priced. Fourth, guarantees are typically reported from the parameters of the mechanism
rather than audited by attack; where an audit is conducted it is performed at
the granularity of the record, and a null result is seldom accompanied by the
evidence that would distinguish an effective defence from an inadequate
adversary. Fifth, and cutting across the foregoing, the cost of a federated
design is reported in the currency of training --- rounds to convergence,
bytes exchanged, energy expended --- while the delay experienced by the
clinician who queries the deployed model is left unmeasured, notwithstanding
that it is this delay which determines whether the system can be used at the
point of care at all.

The present work addresses these five points together. It combines a two-tier
hierarchy with a guarantee whose sensitivity is enforced at the level of the
individual at the point of clipping rather than recovered by conversion
afterwards, and calibrates the noise against every client simultaneously so
that the interposition of the edge tier is reflected in the reported budget. It
measures the utility conceded by moving the unit from the record to the
individual, and characterises that cost as a function of the number of records
each individual contributes rather than at a single operating point. It
enforces and verifies the disjointness of the pretraining population, and
reports the utility forgone in order that the resulting guarantee be valid
end-to-end. Finally, it audits the guarantee by mounting membership inference
at the level of the individual, using both a threshold attack and a
likelihood-ratio attack calibrated per individual, and reports every result
alongside a positive control on a configuration known to memorise, so that a
negative finding may be attributed to the defence rather than to the adversary.
Finally, it specifies and instruments the deployment path, so that the
end-to-end latency of a prediction, the volume of traffic the hierarchy
carries, and the energy each tier expends are reported for the same
configuration that carries the privacy guarantee, rather than for a
differently-configured system.
